\documentclass{article}
\usepackage[utf8]{inputenc}
\usepackage{amsmath, amssymb}
\usepackage{graphicx}
\usepackage{hyperref}
\usepackage{siunitx}
\usepackage{booktabs}
\title{Teori for håndholdt EMP-generator}
\author{}
\date{}

\begin{document}
\maketitle

\section{Innledning}
En elektromagnetisk puls (EMP) er en kort, intens utbrudd av elektromagnetisk energi. I dette prosjektet skapes pulsen ved å utlade en stor kondensator gjennom en spole med svært lav induktans. Den raske endringen i strøm ($di/dt$) genererer et tidsvarierende magnetfelt, som ifølge Faradays lov induserer en spenning i nærliggende ledere.

\section{Grunnleggende fysikk}

\subsection{Energilagring i kondensator}
Energien som er lagret i en kondensator med kapasitans $C$ og spenning $V$ er gitt ved
\begin{equation}
E = \frac{1}{2} C V^2.
\end{equation}

\subsection{Utladning gjennom spole}
Når kondensatoren kobles over en spole med induktans $L$ og ohmsk motstand $R$, dannes en serie RLC-krets. Strømmen $i(t)$ er gitt av differensialligningen
\begin{equation}
L \frac{d^2 i}{dt^2} + R \frac{di}{dt} + \frac{1}{C} i = 0.
\end{equation}
For en underdempet krets ($R < 2\sqrt{L/C}$) er løsningen en dempet sinus:
\begin{equation}
i(t) = \frac{V_0}{\omega_d L} e^{-\alpha t} \sin(\omega_d t),
\end{equation}
der
\begin{align}
\alpha &= \frac{R}{2L}, \\
\omega_0 &= \frac{1}{\sqrt{LC}}, \\
\omega_d &= \sqrt{\omega_0^2 - \alpha^2}.
\end{align}

\subsection{Magnetfelt fra en flat spiralspole}
For en flat, sirkulær spole med $N$ vindinger med radier $r_k$, er magnetfeltet i sentrum ($z=0$) gitt ved summen av bidrag fra hver vinding:
\begin{equation}
B(0) = \sum_{k=1}^{N} \frac{\mu_0 I}{2 r_k}.
\end{equation}
For et punkt på aksen i avstand $z$:
\begin{equation}
B(z) = \sum_{k=1}^{N} \frac{\mu_0 I r_k^2}{2 (r_k^2 + z^2)^{3/2}}.
\end{equation}

\subsection{Indusert spenning}
En liten lukket sløyfe med areal $A$ som gjennomtrenges av et tidsvarierende magnetfelt $B(t)$, får indusert en elektromotorisk spenning (ems) gitt av Faradays lov:
\begin{equation}
\mathcal{E} = -\frac{d\Phi}{dt} = -A \frac{dB}{dt}.
\end{equation}

\section{Kretsens virkemåte}

\subsection{Overordnet blokkskjema}
Enheten består av tre hoveddeler:
\begin{enumerate}
\item \textbf{Ladekrets:} En selvoscillerende flyback-omformer som omdanner 45 V DC fra batteriene til 450 V DC for å lade hovedkondensatoren.
\item \textbf{Energilager og utladningskrets:} En 100 µF kondensator og et trigget gnistgap som kobler kondensatoren til spolen når piezo-tenneren aktiveres.
\item \textbf{Antennespole:} En flat spiralspole som omdanner strømpulsen til et magnetfelt.
\end{enumerate}

\subsection{Ladekrets (selvoscillerende flyback)}

Ladekretsen er basert på en selvoscillerende blocking-/flyback-topologi med
MOSFET-en Q1 og transformatoren T1.

Primærviklingen har 10 vindinger. Den prikkmerkede terminalen er koblet til
$+45\,\text{V}$, mens den andre terminalen går til drain på Q1.

En oppstartsmotstand $R_{\text{start}}$ gir en liten initial gatestrøm slik at
Q1 begynner å lede. Tilbakekoblingsviklingen har 5 vindinger og er koblet med
motsatt polaritet relativt til primærviklingen: den prikkmerkede terminalen
ligger mot C3 og gate, mens den andre terminalen ligger mot $+45\,\text{V}$.

Når primærstrømmen begynner å øke, induseres dermed en spenning i
tilbakekoblingsviklingen som driver gate i positiv retning. Dette gir den
regenerative tilbakekoblingen som er nødvendig for selvoscillasjon.

Gate-spenningen begrenses av gatebeskyttelsen D2/D4, mens R3 gir en definert
utladningsvei for gate. Den nøyaktige turn-off-mekanismen og
svitsjefrekvensen avhenger av transformatorens magnetiske egenskaper,
tilbakekoblingen, gatekretsen og parasittiske komponenter. Disse størrelsene
må derfor verifiseres ved måling og bør ikke betraktes som fast bestemt av
den forenklede modellen alene.

Sekundærviklingen har 200 vindinger og motsatt polaritet relativt til
primæren. Den prikkmerkede sekundærterminalen ligger mot C1-/retur. Under
Q1s ON-periode er likeretteren derfor sperret. Når Q1 slår av, reverseres
viklingsspenningene, sekundærdiodene leder, og energien overføres til C1.

D5, R2 og C2 utgjør en RCD-clamp for å begrense drain-transienter forårsaket
av lekkinduktans. De endelige snubberverdiene må verifiseres ved måling av
$V_{DS}$.

\subsection{Gnistgap og trigger}
Gnistgapet består av to wolframelektroder med en avstand på ~0,5 mm. En piezoelektrisk tenner genererer en høy spenningspuls (>10 kV) som ioniserer gapet, slik at hovedgnisten kan dannes.

\subsection{Utladningsforløp}
Når gnistgapet lukkes, kobles den oppladede kondensatoren direkte over spolen. Kretsen er en underdempet RLC-krets med svært lav motstand. Strømmen stiger til flere tusen ampere i løpet av mikrosekunder, og magnetfeltet når sitt maksimum. Deretter svinger energien mellom kondensatoren og spolen, og pulsen dør ut etter noen få perioder.
\end{document}
